\documentclass[10pt,a4paper]{amsart}
\usepackage[margin=19mm]{geometry}
\usepackage{amsmath,amssymb,mathtools}
\usepackage[scr=rsfs,cal=euler]{mathalfa}
\usepackage{xfrac}
\usepackage[unicode,hidelinks,bookmarks=false]{hyperref}
\newcommand{\ie}{i.e.\ }
\newcommand{\cf}{cf.\ }
\newcommand{\resp}{respectively\ }
\newcommand{\Subsection}{\subsection}
\providecommand{\nobreakdash}{\nobreak\hbox{-}}
\setlength{\emergencystretch}{2em}
\multlinegap0pt
% The published article numbers its sections with Roman numerals; reproduced here.
\renewcommand{\thesection}{\Roman{section}}\let\labelindent\relax
\usepackage{enumitem}
\DeclareMathOperator{\rank}{rank}
\DeclareMathOperator{\tr}{tr}
\newcommand{\set}[2]{\left\{#1 \mid #2\right\}}
\newtheorem{theorem}{Theorem}
\newtheorem{corollary}[theorem]{Corollary}
\newtheorem{lemma}[theorem]{Lemma}
\newtheorem{example}[theorem]{Example}
\newtheorem{assumption}[theorem]{Assumption}
\newtheorem{proposition}[theorem]{Proposition}
\newtheorem{problem}{Problem}
\newtheorem{alg}{Algorithm}
\newtheorem{remark}[theorem]{Remark}
\newtheorem{fact}{Fact}
\newtheorem{definition}[theorem]{Definition}
\newtheorem{notation}{Notation}
\newtheorem{note}{Note}
\newtheorem{claim}{Claim}
\newtheorem{observation}{Observation}
% The publisher class defines the end-of-proof box as a filled 1.3ex square; reproduced here
% because the ieeecolor class file is not present in the runtime and is not redistributed.
\providecommand{\QEDclosed}{\mbox{\rule[0pt]{1.3ex}{1.3ex}}}
\providecommand{\QED}{\QEDclosed}
% Editorial compatibility definition for the publisher's keywords environment.
\newenvironment{IEEEkeywords}{\par\medskip\noindent\textbf{Index Terms}---\ignorespaces}{\par\medskip}
\title{Data-Driven Stabilization Using Prior Knowledge on Stabilizability and Controllability}
\author{Amir Shakouri, Henk J. van Waarde, Tren M.J.T. Baltussen, and W.P.M.H. Heemels}
\date{Source: arXiv:2510.25452v3 (CC BY 4.0); the authors' own native \LaTeX{} source, set here in the editorial AMS \texttt{amsart} wrapper (the IEEE \texttt{ieeecolor} class named by the source is not present in the runtime and is not redistributed).}
\begin{document}
\maketitle
\noindent\textit{Editorial scope.} The counted claim of this unit is the article's Theorem 5 (source label \texttt{th:1}), the theorem stating that for a controllable true system the data are $\Sigma_\text{cont}$--informative for stabilization exactly when they are informative for stabilization, and that a gain stabilizing every data-consistent controllable system stabilizes every data-consistent system. It is delivered together with the authors' complete original run-in proof as the single primary proof body. Retained uncounted support, in source order: the abstract and the keywords block (source0.tex lines 117--123); the whole Problem Formulation section with Definition 1, Proposition 2, Definition 3, Problem 1 and Example 4 (source lines 146--299); the paragraph after the counted statement together with Example 6 (source lines 310--325); and, after the publisher figure, the three auxiliary results with their complete proofs, Lemma 7 with its proof (source lines 335--348), Lemma 8 (source lines 352--355) and Lemma 9 with its proof (source lines 359--365), followed by the notation paragraph that prepares the counted proof (source lines 332--372). Omitted: the Introduction (source lines 125--145), because it refers forward to Theorem 15, Proposition 16 and Section IV-B, which lie outside this source-local core, so that every cross-reference inside the delivered passages resolves inside the delivered document; the publisher figure environment of Figure 1 with its EPS file (source lines 326--331), which is not delivered, and whose only textual mention, the words shown in Fig.~1 inside Example 6, is plain text rather than a cross-reference; the later sections on the rank-deficient case and on the numerical example; and the file preamble of main.tex, replaced by the editorial AMS \texttt{amsart} wrapper declared above. Printed numbering: the authors' own theorem-environment declarations are carried unchanged, including the single shared counter of the theorem, corollary, lemma, example, assumption, proposition, remark and definition environments, so the delivered PDF prints Definition 1, Proposition 2, Definition 3, Example 4, Problem 1, Theorem 5, Example 6, Lemma 7, Lemma 8 and Lemma 9 exactly as the published article does. The wrapper is additionally set to the published Roman section numbering and the section counter is advanced past the omitted Introduction, so the delivered PDF also prints the section headings with the published numbers, II for Problem Formulation and III for Controllability as Prior Knowledge. Class substitution: the source class is the IEEE \texttt{ieeecolor} class together with the local \texttt{generic.sty} stub; neither file is present in the runtime and neither is redistributed. The wrapper is AMS \texttt{amsart} carrying the authors' own operator, macro and theorem declarations, and the publisher's keywords environment and end-of-proof box are reproduced as compatibility definitions that print the keywords line and a filled end-of-proof box exactly as the publisher class does. The publisher's \texttt{showonlyrefs} setting of \texttt{mathtools} is deliberately not carried, so every displayed equation of the retained passages is numbered in the delivered PDF whereas the published article numbers only referenced equations; the printed numbers of the theorem-like environments, which identify the counted claim, are unaffected. Reference handling: the retained passages contain exactly three citation commands, all of them inside the optional titles of Definition 1, Proposition 2 and Lemma 8. The article's bibliography exists only in the authors' bibliography data file \texttt{biblo.bib}; the article's own TeX contains no reference list, and that data file is kept verbatim in the eprint directory and is not redistributed. Each of the three citation commands is therefore delivered as a reference-only adaptation, namely a hyperlink to the article's own abstract page carrying the published citation text [3, Def. 12], [3, Thm. 17] and [27, 2.4.P10] as plain text, so that every citation keeps the number it prints in the published article and the reader can locate the entry in the original article. No reference list is delivered and no reference entry is transcribed or invented. No mathematics is written, repaired or re-derived here: all mathematics is the authors' own native text.
\begin{abstract}
In this work, we study data-driven stabilization of linear time-invariant systems using prior knowledge of system-theoretic properties, specifically stabilizability and controllability. To formalize this, we extend the concept of data informativity by requiring the existence of a controller that stabilizes all systems consistent with the data \emph{and} the prior knowledge. We show that if the system is controllable, then incorporating this as prior knowledge does not relax the conditions required for data-driven stabilization. Remarkably, however, we show that if the system is stabilizable, then using this as prior knowledge leads to necessary and sufficient conditions that are \emph{weaker} than those for data-driven stabilization without prior knowledge. In other words, data-driven stabilization is easier if one knows that the underlying system is stabilizable. We also provide new data-driven control design methods in terms of linear matrix inequalities that complement the conditions for informativity.
\end{abstract}

\begin{IEEEkeywords}
Data-driven control, stabilization, prior knowledge, controllability, stabilizability
\end{IEEEkeywords}

\setcounter{section}{1}
\section{Problem Formulation}
\label{sec:II}

Let $n,m\in\mathbb{N}$. Consider the LTI system
\begin{equation}
\label{eq:1}
x(t+1)=A_\text{true}x(t)+B_\text{true}u(t),
\end{equation}
referred to as the \emph{true system}, where $t\in\mathbb{Z}_+$ denotes time, \mbox{$x(t)\in\mathbb{R}^n$} is the state, and $u(t)\in\mathbb{R}^m$ is the input. The system matrices 
\begin{equation}
(A_\text{true},B_\text{true}) \in \mathcal{M} \coloneqq \mathbb{R}^{n \times n} \times \mathbb{R}^{n \times m}
\end{equation}
are assumed to be unknown. However, we have access to input-state data of the form
\begin{equation}
\label{eq:data}
\mathcal{D}\coloneqq\left(\begin{bmatrix}
    u(0)  & \cdots & u(T-1)
\end{bmatrix}, \begin{bmatrix}
    x(0) & \cdots & x(T)
\end{bmatrix}\right)
\end{equation} 
collected from \eqref{eq:1} within the time horizon $T\in\mathbb{N}$. Given $\mathcal{D}$, we define the matrices \mbox{$X_{-}\coloneqq \begin{bmatrix}
x(0) & x(1) & \cdots & x(T-1)
\end{bmatrix}$} and $X_{+}\coloneqq \begin{bmatrix}
x(1) & x(2) & \cdots & x(T)
\end{bmatrix}$. We also define the matrix \mbox{$U_- = \begin{bmatrix} u(0) & u(1) & \cdots & u(T-1) \end{bmatrix}$}. 

\subsection{Recap of Data-driven stabilization}

Roughly speaking, data-driven stabilization aims at solving the following problem:

\begin{center}
\textit{Given the data $\mathcal{D}$, find a $K$ such that $A_\textup{true}+B_\textup{true}K$ is Schur.}
\end{center}
Based on the collected input-state data, the true system satisfies 
\begin{equation}
\label{eq:lineqtrue}
X_+=A_\textup{true}X_-+B_\textup{true}U_-.
\end{equation}
However, the true system may not be the only one that satisfies this identity. Therefore, a feedback gain that guarantees the stability of the true system must stabilize \emph{all} systems $(A,B)$ satisfying \mbox{$X_+=AX_-+BU_-$}. Therefore, we define the set of \emph{data-consistent systems} as
\begin{equation}
\label{eq:exp}
\Sigma_\mathcal{D}\coloneqq \set{(A,B)\in\mathcal{M}}{X_+=AX_-+BU_-},
\end{equation}
and we sharpen the data-driven stabilization problem as follows:
\begin{center}
\textit{Given the data $\mathcal{D}$, find a $K$ such that $A+BK$ is Schur for all $(A,B)\in\Sigma_\mathcal{D}$.}
\end{center}
The feasibility of this problem depends on the given data. To elaborate on this, we recap the following notion of data informativity. 

\begin{definition}[{\href{https://arxiv.org/abs/2510.25452v3}{[3, Def. 12]}}]
\label{def:1}
The data $\mathcal{D}$ are called \emph{informative for stabilization} if there exists a $K\in\mathbb{R}^{m\times n}$ such that $A+BK$ is Schur for all $(A,B)\in\Sigma_\mathcal{D}$. 
\end{definition}

The following result provides a necessary and sufficient LMI condition for the informativity of the data for stabilization.

\begin{proposition}[{\href{https://arxiv.org/abs/2510.25452v3}{[3, Thm. 17]}}]
\label{prop:1}
The data $\mathcal{D}$ are informative for stabilization \emph{if and only if} there exists $\Theta\in\mathbb{R}^{T\times n}$ such that
\begin{equation}
\label{eq:dd_lmi}
X_-\Theta=\Theta^\top X_-^\top\ \text{ and }\ \begin{bmatrix}
X_-\Theta & X_+\Theta \\
\Theta^\top X_+^\top & X_-\Theta
\end{bmatrix}>0.
\end{equation}
Moreover, $A+BK$ is Schur for all $(A,B)\in\Sigma_\mathcal{D}$ \emph{if and only if} $K=U_-\Theta(X_-\Theta)^{-1}$ for some $\Theta$ satisfying \eqref{eq:dd_lmi}. 
\end{proposition}

Based on Proposition~\ref{prop:1}, a necessary condition for the informativity of the data for stabilization is that $\rank X_-=n$. This condition requires the number of data samples to satisfy $T\geq n$.

\subsection{Data-driven stabilization using prior knowledge}

Let $\Sigma_\text{pk}\subseteq\mathcal{M}$ be a set capturing our prior knowledge of the true system, i.e.,
\begin{equation}
(A_\text{true},B_\text{true})\in\Sigma_\text{pk}.
\end{equation}
Using this prior knowledge, we extend the data-driven stabilization problem to the following.

\begin{center}
\textit{Given the data $\mathcal{D}$ and the set of prior knowledge $\Sigma_\textup{pk}$, find a $K$ such that $A+BK$ is Schur for all $(A,B)\in\Sigma_\mathcal{D}\cap\Sigma_\textup{pk}$.}
\end{center}

The feasibility of this problem depends on the given data \emph{and} the prior knowledge. To study this problem, we extend the notion of data informativity in Definition~\ref{def:1} to the following, which takes the prior knowledge into account. 

\begin{definition}
\label{def:2}
The data $\mathcal{D}$ are called $\Sigma_\textup{pk}$--\emph{informative for stabilization} if there exists a $K\in\mathbb{R}^{m\times n}$ such that $A+BK$ is Schur for all $(A,B)\in\Sigma_\mathcal{D}\cap\Sigma_\text{pk}$. 
\end{definition}

We note that $\mathcal{M}$--informativity for stabilization is equivalent to informativity for stabilization in the sense of Definition~\ref{def:1}.  In this work, we are interested in two important sets of prior knowledge that capture the \emph{stabilizability} and \emph{controllability} of the true system. We denote by $\Sigma_\text{cont}, \Sigma_\text{stab} \subset \mathcal{M}$ the subsets of controllable and stabilizable systems, respectively.
%\begin{equation}
%\begin{split}
%\Sigma_\text{cont}&\coloneqq\set{(A,B)\in\mathcal{M}}{(A,B)\text{ is controllable}},\ \text{and} \\
%\Sigma_\text{stab}&\coloneqq\set{(A,B)\in\mathcal{M}}{(A,B)\text{ is stabilizable}}.
%\end{split}
%\end{equation}

The following example demonstrates that, in case
\begin{equation}
\label{eq:pk=stab}
\Sigma_\text{pk}=\Sigma_\text{stab},
\end{equation}
the conditions for $\Sigma_\text{pk}$-informativity for stabilization are in general \emph{weaker} than those for informativity for stabilization.

\begin{example}
\label{ex:1}
Consider the input data $u(0)=1$, \mbox{$u(1)=2$}, and $u(2)=-1$, and the state data $x(0)=\begin{bmatrix}
1 & 0
\end{bmatrix}^\top$, \mbox{$x(1)=\begin{bmatrix}
2 & 0
\end{bmatrix}^\top$}, $x(2)=\begin{bmatrix}
4 & 0
\end{bmatrix}^\top$, and $x(3)=\begin{bmatrix}
3 & 0
\end{bmatrix}^\top$. The set of data-consistent systems reads
\begin{equation}
\Sigma_\mathcal{D}=\set{\left(\begin{bmatrix}
1 & \alpha \\
0 & \beta
\end{bmatrix},\begin{bmatrix}
1 \\ 0
\end{bmatrix}\right)}{\alpha,\beta\in\mathbb{R}}.
\end{equation}
It follows from Proposition~\ref{prop:1} that the data are not informative for stabilization since $X_-=\begin{bmatrix}
1 & 2 & 4 \\
0 & 0 & 0
\end{bmatrix}$ does not have full row rank. However, we have
\begin{equation}
\Sigma_\mathcal{D}\cap\Sigma_\text{stab}=\set{\left(\begin{bmatrix}
1 & \alpha \\
0 & \beta
\end{bmatrix},\begin{bmatrix}
1 \\ 0
\end{bmatrix}\right)}{\alpha\in\mathbb{R},|\beta|<1}.
\end{equation}
It is evident that $K=\begin{bmatrix}
-1 & 0
\end{bmatrix}$ is a stabilizing feedback gain for all the systems in \mbox{$\Sigma_\mathcal{D}\cap\Sigma_\text{stab}$}. Therefore, the data are \mbox{$\Sigma_\text{stab}$--informative} for stabilization. 
\end{example}

Example~\ref{ex:1} shows that data-driven stabilization using stabilizability as prior knowledge may be possible in case Proposition~\ref{prop:1} fails to provide a feedback gain for all data-consistent systems. This motivates the study of stabilizability as prior knowledge for data-driven stabilization. Another closely related prior knowledge that is studied in this paper is controllability. Formally, we thus consider the following problem. 
\begin{problem}
\label{prob:1}
Find necessary and sufficient conditions under which the data $\mathcal{D}$ are (i) $\Sigma_\text{cont}$--informative for stabilization; \mbox{(ii) $\Sigma_\text{stab}$--informative} for stabilization.
\end{problem}



\section{Controllability as Prior Knowledge}

The main result of this section is the following theorem presenting the solution for Problem~\ref{prob:1}(i). This theorem shows that controllability as prior knowledge is \emph{not} useful, i.e., the data-driven stabilization using this prior knowledge is equivalent to data-driven stabilization without any prior knowledge. 


\begin{theorem}
\label{th:1}
Suppose that $(A_\text{true},B_\text{true})\in\Sigma_\text{cont}$. Then, the following statements are equivalent:
\begin{enumerate}[label=(\alph*),ref=\ref{th:1}(\alph*)]
    \item The data $\mathcal{D}$ are $\Sigma_\text{cont}$--informative for stabilization.
    \item The data $\mathcal{D}$ are informative for stabilization.
\end{enumerate}
Moreover, if $K$ is such that $A+BK$ is Schur for all \mbox{$(A,B)\in\Sigma_\mathcal{D}\cap\Sigma_\text{cont}$}, then $A+BK$ is Schur for all $(A,B)\in\Sigma_\mathcal{D}$. 
\end{theorem}


The following example provides an intuition on the fact that controllability as prior knowledge does not affect the conditions required for data-driven stabilization. 

\begin{example}
\label{ex:2}
For the sake of illustration, consider a system with $n=1$ and $m=2$ as follows:
\begin{equation}
x(t+1)=ax(t)+\begin{bmatrix}
b_1 & b_2
\end{bmatrix}u(t).
\end{equation}
Starting from $x(0)=-1$, we apply $u(0)=\begin{bmatrix}
1 & -1
\end{bmatrix}^\top$ and we measure $x(1)=-1$. Given these data, the set $\Sigma_\mathcal{D}$ consists of all systems with parameters $a$, $b_1$, and $b_2$ lying on the plane shown in Fig. 1. The only uncontrollable system on this plane corresponds to $a=1$ and $b_1=b_2=0$, which is shown by the red dot. Now, Theorem~\ref{th:1} states that if a feedback gain stabilizes all systems on the plane excluding the one shown in red, then it also stabilizes the red point. 
\end{example}



To prove Theorem~\ref{th:1}, we need three auxiliary results. The first result is the following lemma, showing that if a parameterized family of systems is controllable at a single point, then it is controllable at almost all points. 

\begin{lemma}
\label{lem:cont_direct}
Let $(M,N)\in\Sigma_{\text{cont}}$, $M_0\in\mathbb{R}^{n\times n}$, and $N_0\in\mathbb{R}^{n\times m}$. Then, the pair $(M+\alpha M_0,N+\alpha N_0)$ is controllable for all but at most $n^2$ values of $\alpha\in\mathbb{R}$.

\end{lemma}
\begin{proof}
Denote the Kalman controllability matrix of the pair $(M+\alpha M_0,N+\alpha N_0)$ by
\begin{equation}
\mathcal{C}(\alpha)\coloneqq\begin{bmatrix}
N+\alpha N_0 & \cdots & (M+\alpha M_0)^{n-1}(N+\alpha N_0)
\end{bmatrix}.
\end{equation}
Let $\bar{\mathcal{C}}(\alpha)\in\mathbb{R}^{n\times n}$ be a square submatrix of $\mathcal{C}(\alpha)$ such that $\bar{\mathcal{C}}(0)$ is nonsingular. Observe that the entries of $\bar{\mathcal{C}}(\alpha)$ are polynomials of $\alpha$ of degree at most $n$. It can be concluded from, e.g., Leibniz's formula for determinants, that $\det(\bar{\mathcal{C}}(\alpha))$ is a polynomial of $\alpha$ of degree at most $n^2$. This polynomial has at most $n^2$ roots counting multiplicity, which implies that $\rank \bar{\mathcal{C}}(\alpha)<n$ for at most $n^2$ distinct values \mbox{of $\alpha$}. Therefore, $\mathcal{C}(\alpha)$ is rank deficient for at most $n^2$ distinct values of $\alpha$.
\end{proof}

The second auxiliary result is the following lemma, providing a necessary and sufficient condition for two matrices to share the same eigenvalues.

\begin{lemma}[{\href{https://arxiv.org/abs/2510.25452v3}{[27, 2.4.P10]}}]
\label{lem:horn_2.4.P10}
Let $M,N\in\mathbb{R}^{n\times n}$. Then, $M$ and $N$ share the same eigenvalues with the same algebraic multiplicities \emph{if and only if} for every $k\in\{1,\ldots,n\}$ we have $\tr(M^k)=\tr(N^k)$. 
\end{lemma}

The third auxiliary result discusses the stability of matrix pencils, which is presented in the following lemma.

\begin{lemma}
\label{lem:A+delta*B}
Let $\varepsilon\in\mathbb{R}$, $\mathcal{F}\subset\mathbb{R}$ be a finite set, and $M,N\in\mathbb{R}^{n\times n}$ be such that $M+\delta N$ is Schur for all $\delta\in[\varepsilon,\infty)\backslash\mathcal{F}$. Then, $N$ is nilpotent and $M+\delta N$ is Schur for all $\delta\in\mathbb{R}$. 
\end{lemma}
\begin{proof}
Let $\bar{\delta}>0$ be sufficiently large such that $[\bar{\delta},\infty)\cap\mathcal{F}=\varnothing$. We note that the trace of a matrix is the sum of its eigenvalues. Since $M+\delta N$ is Schur for all $\delta\in[\bar{\delta},\infty)$, we have $|\tr ((M+\delta N)^k)|\leq n$ for all $\delta\in[\bar{\delta},\infty)$ and all $k\in\mathbb{N}$. Note that for every $k\in\mathbb{N}$, \mbox{$p_k(\delta)= \tr ((M+\delta N)^k)$} is a polynomial of degree at most $k$. The boundedness of $p_k(\delta)$ in the interval $[\bar{\delta},\infty)$ implies that $p_k(\delta)$ is constant, i.e., the coefficients of $\delta^i$ are zero for all $i\in\{1,\ldots,k\}$. In particular, the coefficient corresponding to $\delta^k$, i.e., $\tr(N^k)$, is equal to zero. Thus, $\tr(N^k)=0$ for all $k\in\mathbb{N}$. Based on Lemma~\ref{lem:horn_2.4.P10}, this implies that all eigenvalues of $N$ are equal to zero. Thus, $N$ is nilpotent. In addition, since $p_k(\delta)$ is constant for all $k\in\mathbb{N}$, we have $p_k(\delta)=p_k(0)$, thus, \mbox{$\tr(M^k)=\tr((M+\delta N)^k)$} for all $k\in\mathbb{N}$. Now, we use Lemma~\ref{lem:horn_2.4.P10} to conclude that $M$ and $M+\delta N$ share the same eigenvalues for all $\delta\in\mathbb{R}$. Therefore, $M+\delta N$ is Schur for all $\delta\in\mathbb{R}$. 
\end{proof}

The proof of Theorem~\ref{th:1} now follows from Lemmas~\ref{lem:cont_direct} and~\ref{lem:A+delta*B}. To facilitate the proof, we introduce the following notation:
\begin{equation}
\Sigma_\mathcal{D}^0\coloneqq\set{(A_0,B_0)\in\mathcal{M}}{A_0X_-+B_0U_-=0}.
\end{equation}
It is evident that the set $\Sigma_\mathcal{D}^0$ satisfies $\Sigma_\mathcal{D}^0+\Sigma_\mathcal{D}=\Sigma_\mathcal{D}$.


\textit{Proof of Theorem~\ref{th:1}:} 
(b)$\Rightarrow$(a): This implication is evident since if $K\in\mathbb{R}^{m\times n}$ is such that $A+BK$ is Schur for all $(A,B)\in\Sigma_\mathcal{D}$, then $A+BK$ is Schur for all $(A,B)\in\Sigma_\mathcal{D}\cap\Sigma_\text{cont}$.  

(a)$\Rightarrow$(b): This implication obviously holds if $\Sigma_\mathcal{D}\subseteq\Sigma_\text{cont}$. Now, assume that there exists $(\bar{A},\bar{B})\in\Sigma_\mathcal{D}$ that is not controllable. Let $K$ be such that $A+BK$ is Schur for all $(A,B)\in\Sigma_\mathcal{D}\cap\Sigma_\text{cont}$. It suffices to show that $\bar{A}+\bar{B}K$ is Schur. To that end, note that the pair $(\bar{A},\bar{B})$ satisfies $X_+=\bar{A}X_-+\bar{B}U_-$. Since the true system also satisfies $X_+=A_\textup{true}X_-+B_\textup{true}U_-$, we have
$\bar{A}=A_\text{true}+A_0$ and $\bar{B}=B_\text{true}+B_0$
for some \mbox{$(A_0,B_0)\in\Sigma_\mathcal{D}^0$}. We note that \mbox{$(A_\text{true}+\alpha A_0,B_\text{true}+\alpha B_0)\in\Sigma_\mathcal{D}$} for all $\alpha\in\mathbb{R}$. Based on Lemma~\ref{lem:cont_direct}, since $(A_\text{true},B_\text{true})\in\Sigma_\text{cont}$, we have \linebreak $(A_\text{true}+\alpha A_0,B_\text{true}+\alpha B_0)\in\Sigma_\mathcal{D}\cap\Sigma_\text{cont}$ for all but a finite number of $\alpha \in \mathbb{R}$. Thus, $A_\text{true}+B_\text{true}K+\alpha(A_0+B_0 K)$ is Schur for all but a finite number of $\alpha \in \mathbb{R}$. Based on Lemma~\ref{lem:A+delta*B}, this implies that $A_\text{true}+B_\text{true}K+\alpha(A_0+B_0 K)$ is Schur for all $\alpha\in\mathbb{R}$. We take $\alpha=1$ and conclude that $\bar{A}+\bar{B}K$ is Schur. \hfill \QED

\end{document}
